\section{Conclusion}
\label{sec:conclusion}

This study examined commodity tail dependence and exchange-rate risk in Ghana using two co-equal calendar constructions, explicit marginal-model adequacy checks, eight static bivariate copula families, finite-threshold tail-concentration measures and multiplicity-controlled stress inference. The purpose was not to identify a single preferred preprocessing rule or a single best-fitting copula, but to determine which conclusions remain defensible after calendar alignment, marginal specification, absolute model adequacy and multiple testing are treated explicitly.

Three findings survive that scrutiny. First, the stress-period conclusion is calendar-robust. Across the combined COVID-19 plus 2024 definition, COVID-19 alone and 2024 alone, neither the business-day panel nor the common-observation panel produces a Bonferroni- or Benjamini-Hochberg-adjusted discovery among the 60 finite-tail comparisons in each family. This result does not establish that dependence was unchanged during those episodes. It means that the tested stress-versus-calm differences do not provide multiplicity-adjusted evidence strong enough to support individual changes under the prespecified thresholds and bootstrap design.

Second, absolute copula adequacy is itself sensitive to the way the observation calendar is constructed. Gold-Brent, Gold-WTI and Brent-WTI reject all eight tested static families in both panels, giving a calendar-robust indication that the candidate set is inadequate for these pairs. The same conclusion does not hold for Cocoa-Gold, Cocoa-Brent and Cocoa-WTI. Each rejects all eight families under the business-day construction but only a subset under common-observation alignment. These cocoa-related results show that calendar choice can alter higher-order dependence conclusions even when the corresponding return series remain strongly correlated across panels. For applied dependence analysis, calendar construction is therefore part of the statistical specification rather than a neutral preprocessing detail.

Third, the Cedi remains the main unresolved marginal problem. None of the tested continuous marginal candidates passes the full adequacy gate under either calendar construction. Commodity-Cedi copula non-rejections are stable across panels, but they should not be interpreted as proof that the tested families are correct. Likewise, the Cedi EVT results indicate a heavy fitted loss tail under the retained reference model, but those quantities remain conditional on marginal misspecification. The separate clipping audit shows that the identified Panel A PIT tie is numerically minor and does not drive the substantive conclusions. The larger issue is the inability of the tested continuous marginal models to reproduce the zero-heavy daily Cedi process adequately.

The paper's contribution is therefore methodological as much as empirical. It shows how apparently simple conclusions about commodity dependence can change once absolute fit, calendar construction and multiple testing are considered jointly. It also demonstrates the value of retaining negative and inconvenient findings rather than forcing a single narrative. Some conclusions survive both panels, some become calendar-sensitive, one historical analysis branch is withdrawn after its data fail audit, and Cedi-related evidence remains explicitly conditional. This produces a narrower set of claims than the earlier manuscript, but a more defensible one.

For Ghana, the results should not be read as evidence that commodity markets are unimportant for the Cedi or that crisis-period dependence cannot increase. The analysis addresses a specific daily-source-frequency tail question and does not estimate causal pass-through, fiscal effects, reserve implications or lower-frequency macroeconomic transmission. Those questions remain open. Future work would benefit from a separately validated macroeconomic data pipeline, more flexible models for the Cedi's mixed discrete-continuous behavior, and dynamic, regime-switching, vine or nonparametric dependence models for commodity pairs that fail the static adequacy tests.

The broader lesson is straightforward. In multivariate risk analysis, robustness is not demonstrated by one statistically significant cell, one preferred calendar or one lowest-AIC model. The most credible conclusions are those that survive alternative defensible data constructions, explicit model-adequacy checks and the inferential burden created by multiple comparisons.
